% we include the glossary here (frontmatter is included with \input, so this command is as if it was in main.tex)
\newglossaryentry{analiseemocional}{
 name={análise emocional},
 description={Processo computacional de identificação de emoções expressas linguisticamente em texto, com base em padrões lexicais, semânticos e discursivos, sem inferência directa sobre estados psicológicos internos do utilizador.}
}

\newglossaryentry{emocaoexpressa}{
 name={emoção expressa},
 description={Manifestação linguística de um estado afectivo através do discurso, distinta da emoção sentida, e constituindo a única dimensão emocional directamente observável e passível de modelação computacional.}
}

\newglossaryentry{ambivalenciaemocional}{
 name={ambivalência emocional},
 description={Coexistência, no mesmo enunciado, de pelo menos uma emoção de valência positiva e uma negativa, ambas acima do limiar de activação. É uma propriedade derivada, não uma emoção autónoma. Distingue-se das emoções mistas da literatura, que podem incluir coocorrência sem oposição de valência: duas emoções positivas não contam.}
}

\newglossaryentry{valenciaemocional}{
 name={valência emocional},
 description={Dimensão afectiva que classifica emoções como positivas ou negativas, utilizada para distinguir emoções de natureza oposta no cálculo da ambivalência emocional.}
}

\newglossaryentry{empatiacomputacional}{
 name={empatia computacional},
 description={Capacidade funcional de um sistema artificial gerar respostas que são percebidas como empáticas pelo utilizador, avaliada com base na adequação emocional da resposta e não na existência de estados emocionais internos no sistema.}
}

\newglossaryentry{empatiapercebida}{
 name={empatia percebida},
 description={Avaliação subjectiva realizada por utilizadores ou avaliadores humanos relativamente ao grau de compreensão, sensibilidade emocional e adequação demonstrados por uma resposta gerada por um sistema conversacional.}
}

\newglossaryentry{classificacaomultilabel}{
 name={classificação multi-label},
 description={Abordagem de classificação em que um mesmo texto pode ser associado simultaneamente a múltiplas categorias emocionais, permitindo a representação da coocorrência de emoções.}
}

\newglossaryentry{propriedadederivada}{
 name={propriedade derivada},
 description={Característica calculada a partir da combinação de saídas de um modelo, não correspondendo a um rótulo supervisionado, mas a uma relação funcional entre variáveis observáveis.}
}

\newglossaryentry{limiaractivacao}{
 name={limiar de activação},
 description={Valor mínimo de activação emocional acima do qual uma emoção é considerada relevante para efeitos de análise e cálculo da ambivalência emocional.}
}

\newglossaryentry{indiceambivalencia}{
 name={índice de ambivalência},
 description={Métrica contínua que representa a intensidade da ambivalência emocional, calculada como função das activa\-ções das emoções positiva e negativa dominantes num enunciado.}
}

\newglossaryentry{pipeline}{
 name={pipeline de processamento},
 description={Artefacto de código que analisa o enunciado, deriva a ambivalência e gera a resposta. Corre nas duas condições experimentais; o que muda entre elas é a presença ou ausência do sinal contextual no \textit{prompt}.}
}

\newglossaryentry{sinalcontextual}{
 name={sinal contextual},
 description={Neste trabalho, o par de emoções dominantes (positiva e negativa) injectado no \textit{prompt} do gerador, sem constituir um rótulo supervisionado nem uma regra fixa.}
}

\newglossaryentry{within}{
 name={within-subject},
 description={Desenho experimental em que os mesmos participantes avaliam todas as condições experimentais, permitindo controlar variabilidade interindividual.}
}

\newglossaryentry{likert}{
 name={Likert},
 description={Instrumento de medição baseado em escalas ordinais, utilizado para recolher avaliações subjectivas de participantes relativamente a dimensões como empatia, adequação emocional e utilidade.}
}

\newglossaryentry{geracaocondicionada}{
 name={geração condicionada},
 description={Processo de geração de texto em que a resposta produzida por um modelo de linguagem é influenciada por instruções ou contexto adicional, sem recurso a re-treinamento do modelo.}
}

\newglossaryentry{inferenciaspsicologicas}{
 name={inferências psicológicas internas},
 plural={inferências psicológicas internas},
 description={Suposições sobre estados mentais ou emocionais internos do utilizador que não são directamente observáveis no texto e que são explicitamente evitadas na abordagem proposta.}
}

\newglossaryentry{texto}{
 name={texto},
 plural={textos},
 description={Linguagem escrita enquanto objecto de processamento computacional (input, conjunto de textos ou output gerado), sem comprometer uma unidade conversacional.}
}

\newglossaryentry{enunciado}{
 name={enunciado},
 plural={enunciados},
 description={Unidade de fala ou escrita produzida por um interlocutor numa dada ocasião. Neste trabalho, é a unidade operacional da ambivalência e da geração: cada estímulo é um enunciado do utilizador, eventualmente composto por mais do que uma frase, classificado de forma isolada.}
}

\newglossaryentry{discurso}{
 name={discurso},
 plural={discursos},
 description={Uso da linguagem em contexto comunicativo. Designa a emoção e a avaliação na linguagem, e não no estado interno, sem coincidir necessariamente com um diálogo completo.}
}

\newglossaryentry{dialogo}{
 name={diálogo},
 plural={diálogos},
 description={Sequência de enunciados entre interlocutores. O histórico do diálogo são os enunciados anteriores ao enunciado corrente.}
}

\newglossaryentry{estimulo}{
 name={estímulo},
 plural={estímulos},
 description={Enunciado construído para o protocolo empírico. Neste trabalho há 12 estímulos em inglês, redigidos pelo autor, filtrados pelo GoEmotions, e cruzados com época e condição.}
}

\newglossaryentry{marcohistorico}{
 name={marco histórico},
 plural={marcos históricos},
 description={Momento da evolução da geração conversacional que justifica escolher um tipo de modelo. Neste trabalho há quatro: \(\sim\)2019--20 (diálogo clássico), \(\sim\)2020--21 (diálogo social), \(\sim\)2023--24 (LLM \textit{instruction-tuned}) e \(\sim\)2024--25 (LLM alinhado recente).}
}

\newglossaryentry{epoca}{
 name={época},
 plural={épocas},
 sort={epoca},
 description={Factor experimental que alinha um marco histórico a um gerador âncora (E1--E4). Não isola o tempo como causa: mistura arquitectura, treino, alinhamento, formato de \textit{prompt} e descodificação.}
}

\newglossaryentry{geradorancora}{
 name={gerador âncora},
 plural={geradores âncora},
 description={Modelo público que ocupa uma época: E1 DialoGPT-medium, E2 BlenderBot, E3 Phi-3-mini-Instruct e E4 Qwen2.5-3B-Instruct.}
}

\newglossaryentry{baseline}{
 name={condição baseline},
 sort={baseline},
 description={Nível do factor experimental em que o gerador âncora produz a resposta sem metadados de ambivalência.}
}

\newglossaryentry{condicaopipeline}{
 name={condição com sinal},
 sort={condicao com sinal},
 description={Nível do factor experimental em que o gerador âncora recebe o par de emoções dominantes positivas e negativas, com a instrução de o usar só como informação de fundo. Nos dados e no CSV, este nível está codificado como \texttt{pipeline}.}
}

\newglossaryentry{celula}{
 name={célula},
 plural={células},
 sort={celula},
 description={Unidade do factorial: um estímulo \(\times\) uma época \(\times\) uma condição (96 no total). No questionário, cada célula apresentada é um item; o juízo Likert de um participante sobre um item é uma classificação.}
}

\newglossaryentry{delta}{
 name={diferencial de empatia},
 sort={delta},
 description={Contraste \(\Delta\) da condição com sinal face à baseline na mesma época (nos dados, \texttt{pipeline} \(-\) \texttt{baseline}). Não mede a empatia absoluta do gerador, mas o ganho (ou perda) associado ao sinal contextual.}
}

\newglossaryentry{filtroestrito}{
 name={filtro estrito},
 sort={filtro estrito},
 description={Recorte da análise principal: sessão concluída, verificação de atenção correcta e sem ritmo demasiado rápido. Nesta amostra, 52 participantes e 2016 classificações.}
}

\newacronym{acl}{ACL}{Association for Computational Linguistics}
\newacronym[shortplural={APIs},longplural={Application Programming Interfaces}]{api}{API}{Application Programming Interface}
\newacronym{bert}{BERT}{Bidirectional Encoder Representations from Transformers}
\newacronym{cli}{CLI}{Command-Line Interface}
\newacronym{cpu}{CPU}{Central Processing Unit}
\newacronym{csv}{CSV}{Comma-Separated Values}
\newacronym{dam}{DAM}{Desvio Absoluto Mediano}
\newacronym{dsr}{DSR}{Design Science Research}
\newacronym{ecm}{ECM}{Emotional Chatting Machine}
\newacronym[shortplural={FKs},longplural={Foreign Keys}]{fk}{FK}{Foreign Key}
\newacronym{hmm}{HMM}{Hidden Markov Model}
\newacronym{ic}{IC}{Intervalo de Confiança}
\newacronym{ieee}{IEEE}{Institute of Electrical and Electronics Engineers}
\newacronym{isep}{ISEP}{Instituto Superior de Engenharia do Porto}
\newacronym{json}{JSON}{JavaScript Object Notation}
\newacronym[shortplural={LLMs},longplural={Large Language Models}]{llm}{LLM}{Large Language Model}
\newacronym{mege}{MEGE}{Mixed Emotion Graph Model for Empathetic Dialogue Generation}
\newacronym{nec}{NEC}{Non-Emotion-Centric}
\newacronym{pal}{PAL}{Persona-Augmented Emotional Support Conversation}
\newacronym{php}{PHP}{Hypertext Preprocessor}
\newacronym[shortplural={PKs},longplural={Primary Keys}]{pk}{PK}{Primary Key}
\newacronym{pln}{PLN}{Processamento de Linguagem Natural}
\newacronym{qa}{QA}{Quality Assurance}
\newacronym{reg}{REG}{Retrieval via Emotion Similarity for Guiding Empathetic Dialogue Generation}
\newacronym{reml}{REML}{Restricted Maximum Likelihood}
\newacronym{se}{SE}{Standard Error}
\newacronym{sql}{SQL}{Structured Query Language}
\newacronym{strideed}{STRIDE-ED}{Strategy-Grounded Stepwise Reasoning Framework for Empathetic Dialogue Systems}
\newacronym{tmdei}{TMDEI}{Mestrado em Engenharia Informática}


\frontmatter % Use roman page numbering style (i, ii, iii, iv...) for the pre-content pages

\pagestyle{plain} % Default to the plain heading style until the thesis style is called for the body content

%----------------------------------------------------------------------------------------
%	OFFICIAL ISEP COVER (front A4) + TITLE PAGE
%----------------------------------------------------------------------------------------
% The portal export is a hardcover wrap (back + spine + front). Published
% dissertations use only the front face, A4, as page 1; the LaTeX title page
% follows immediately as page 2.

\includepdf[
	pages=1,
	pagecommand={\thispagestyle{empty}\ifdef{\pdfbookmark}{\pdfbookmark[0]{Capa}{bookmark-capa-oficial}}{}}
]{frontmatter/assets/capa-oficial-isep-frente-A4.pdf}

% The titlepage environment does \cleardoublepage\newpage, which would insert a
% blank leaf after the official cover. Write the title on the page already opened.
\makeatletter
\let\tmdeinewpage\newpage
\let\tmdeicleardoublepage\cleardoublepage
\let\newpage\relax
\let\cleardoublepage\relax
\maketitlepage
\let\newpage\tmdeinewpage
\let\cleardoublepage\tmdeicleardoublepage
\clearpage
\makeatother


%----------------------------------------------------------------------------------------
%	STATEMENT of INTEGRITY
%----------------------------------------------------------------------------------------
\integritystatement

%----------------------------------------------------------------------------------------
%	DEDICATION  (optional)
%----------------------------------------------------------------------------------------

\begin{dedicatory}
\begin{center}
\textit{Aos que nunca deixaram de perguntar porquê.}
\end{center}
\end{dedicatory}

%----------------------------------------------------------------------------------------
%	ABSTRACT PAGE
%----------------------------------------------------------------------------------------

\begin{abstract}

Os sistemas conversacionais baseados em Processamento de Linguagem Natural têm vindo a ganhar espaço na interacção humano-máquina. Mesmo assim, a análise emocional em texto continua limitada perante emoções complexas e coexistentes. Este trabalho centra-se na ambivalência emocional, a coexistência de emoções de valências distintas no mesmo enunciado, e no problema de saber se um sinal explícito dessa ambivalência ainda acrescenta empatia percebida em geradores de épocas diferentes.

O objectivo é operacionalizar a ambivalência em texto e medir o seu contributo, face a uma \textit{baseline} sem esse sinal, em quatro geradores âncora associados às épocas E1--E4: DialoGPT, BlenderBot, Phi-3-mini-Instruct e Qwen2.5-3B-Instruct. A ambivalência é definida a partir de um classificador multilabel GoEmotions e injectada como contexto numa pipeline em Python, com Hugging Face Transformers. A avaliação é humana e cega: 52 participantes no filtro estrito, 2016 classificações Likert, num questionário online desenvolvido para o estudo. A inferência principal é um modelo misto cruzado.

A hipótese de um ganho maior nos geradores mais antigos (DialoGPT e BlenderBot), a cair até zero nos mais recentes, não se confirma. Todos os intervalos de confiança do diferencial \(\Delta\) incluem zero. O padrão mais estável é a diferença entre geradores: a empatia percebida varia sobretudo entre geradores âncora (cerca de \(1{,}5\) em DialoGPT, \(6{,}1\) em Phi-3), não entre a condição baseline e a condição com sinal. As diferenças absolutas misturam arquitectura, alinhamento e formato de \textit{prompt}, e o rótulo histórico de cada modelo não isola um factor causal. A explicitação de rótulos afectivos derivados do próprio enunciado não apresentou benefício incremental detectável neste desenho. Isso não prova que informação afectiva seja inútil. A avaliação mede a empatia das respostas, não a ambivalência detectada pelo classificador.

\end{abstract}

\begin{abstractotherlanguage}

Conversational systems based on Natural Language Processing are increasingly used in human-machine interaction. Even so, emotion analysis in text remains limited when affect is complex or mixed. This work focuses on emotional ambivalence, the co-occurrence of differently valenced emotions in the same utterance, and on whether an explicit signal of that ambivalence still adds perceived empathy across generators from different eras.

The aim is to operationalise ambivalence in text and to measure its contribution, against a baseline without that signal, in four anchor generators associated with epochs E1--E4: DialoGPT, BlenderBot, Phi-3-mini-Instruct and Qwen2.5-3B-Instruct. Ambivalence is defined from a GoEmotions multi-label classifier and injected as context in a Python generation pipeline that uses Hugging Face Transformers. The evaluation was conducted by human raters under blinded conditions: 52 participants in the strict-filter sample provided 2,016 Likert ratings through an online questionnaire developed for the study. The primary inferential analysis used a mixed-effects model with crossed random effects.

The hypothesis of a larger gain on the older generators (DialoGPT and BlenderBot), falling towards zero on the more recent ones, is not confirmed. All confidence intervals for the differential \(\Delta\) include zero. The most stable pattern is the difference among generators: perceived empathy varies mainly across anchor generators (about \(1.5\) on DialoGPT, \(6.1\) on Phi-3), not between the baseline condition and the signal condition. Absolute differences mix architecture, alignment and prompt format, and the historical label of each model does not isolate a causal factor. Making affective labels derived from the utterance itself explicit did not yield a detectable incremental benefit in this design. That does not prove that affective information is useless. The evaluation measures the empathy of the replies, not the ambivalence detected by the classifier.

\end{abstractotherlanguage}

%----------------------------------------------------------------------------------------
%	ACKNOWLEDGEMENTS (optional)
%----------------------------------------------------------------------------------------

\begin{acknowledgements}

Agradeço ao Professor Doutor Constantino Martins, pela orientação, pela disponibilidade e pelas críticas que ajudaram a clarificar o desenho e o foco deste trabalho.

Agradeço ao Departamento de Engenharia Informática do ISEP e ao \gls{tmdei} pelo enquadramento académico que tornou possível esta dissertação.

Agradeço a todas as pessoas que participaram na avaliação humana e a quem ajudou a divulgar o questionário. Sem essa colaboração, a parte empírica não avançaria.

Agradeço ainda à minha família e aos amigos próximos, pelo apoio ao longo do percurso e pela paciência nos períodos mais intensos de trabalho.

\end{acknowledgements}

%----------------------------------------------------------------------------------------
%	LIST OF CONTENTS/FIGURES/TABLES PAGES
%----------------------------------------------------------------------------------------

\tableofcontents % Prints the main table of contents

\listoffigures % Prints the list of figures

\listoftables % Prints the list of tables

\iflanguage{portuguese}{
\renewcommand{\listalgorithmname}{Lista de Algoritmos}
}
\cleardoublepage
\chapter*{\listalgorithmname}
\addcontentsline{toc}{chapter}{\listalgorithmname}
\makeatletter
\@starttoc{loa}
\makeatother


%\renewcommand{\lstlistlistingname}{List of Source Code}
%\iflanguage{portuguese}{
%\renewcommand{\lstlistlistingname}{Lista de C\'odigo}
%}
%\lstlistoflistings % Prints the list of listings (programming language source code)
%\addchaptertocentry{\lstlistlistingname}


%----------------------------------------------------------------------------------------
%	SYMBOLS
%----------------------------------------------------------------------------------------

%\begin{symbols}{lll} % Include a list of Symbols (a three column table)

%$a$ & distance & \si{\meter} \\
%$P$ & power & \si{\watt} (\si{\joule\per\second}) \\
%Symbol & Name & Unit \\

%\addlinespace % Gap to separate the Roman symbols from the Greek

%$\omega$ & angular frequency & \si{\radian} \\

%\end{symbols}

%----------------------------------------------------------------------------------------
%	ACRONYMS
%----------------------------------------------------------------------------------------

\newcommand{\listacronymname}{List of Abbreviations}
\iflanguage{portuguese}{
\renewcommand{\listacronymname}{Lista de Siglas e Acrónimos}
}

\printglossary[title=\listacronymname,type=\acronymtype,style=long]
\glsresetall

%----------------------------------------------------------------------------------------
%	DONE
%----------------------------------------------------------------------------------------

\mainmatter % Begin numeric (1,2,3...) page numbering
\pagestyle{thesis} % Return the page headers back to the "thesis" style
